% The v1 study is thesis-only. Numbers come from results/ablation_val and results/ablation_test (tools/aggregate_ablation.py, tables written by
% tools/paper_extras.py), from DECISIONS.md and DEVIATIONS.md (2026-10-27), and from results/final/stats_held.json.
\section{Question and analysis plan}
The first study asked whether a suppressor trained only on CT renders of the public LIDC-IDRI collection keeps lesions, and whether showing
it lesions during training helps. The analysis plan (SAP.md) was written and frozen with \path{freeze.json} on 27 October 2026, before any
locked set was evaluated. Its first hypothesis is synthetic, with exact truth:
\begin{description}
\item[H1] On synthetic test CTs, a lesion-aware arm absorbs less of an inserted lesion than the recipe arm (per-CT mean of $A$, paired,
 patient-cluster bootstrap interval, Holm over the lesion-aware arms), and its bone-image MAE is non-inferior within a margin of 5\,\% of the
 recipe's mean MAE (one-sided 5\,\%).
\end{description}
The same rule, applied to synthetic \emph{validation} CTs, selected the v1 model: an arm qualifies only if its absorption is lower with Holm
$p<0.05$ and its bone MAE is non-inferior. SAP.md locked the real sets: SZCH-X-Rays (75 dual-energy pairs), JSRT, Node21, the TBX11K
validation set and VinDr-RibCXR were never to be used for training, tuning or selection; the validation parts of NIH ChestX-ray14 and
VinBigData could inform pilots. Real-image results could never choose a model.

\section{Design}
Each arm changes one factor of the recipe of Chapter~\ref{ch:data} (the training recipe of the reference paper, re-implemented): the
lesion-aware arms add counterfactual lesion insertion (CLIA, up to three lesions from training CTs pasted into the composite), a loss term
weighted on lesion pixels, or both; the factor arms remove the frequency loss, the histogram-matching augmentation, the marking-ROI loss or
the rotation augmentation. Every arm was trained with three seeds for \num{200000} samples, at the learning rate fixed by a pilot on
synthetic validation CTs ($2\times10^{-4}$; Appendix~\ref{app:deviations}). CTs were split by patient (70/10/20) with a fixed hash; 844 CTs met
the usable-CT rules, and the selection used 98 validation CTs with 177 insertion sites, the test 158 CTs with 317 sites. Lesions for the audit
come from other CTs of the same split, never from training CTs.

\section{Outcome of the selection rule}
On the validation CTs (Table~\ref{tab:v1_ablation}) every lesion-aware arm absorbed less than the recipe arm (CLIA by 0.033, the
lesion-weighted loss by 0.036, the combination by 0.100), but only CLIA kept its bone MAE within the margin: the lesion-weighted loss sat on
the margin ($-0.00036$ against $-0.00036$) and the combination was about 6\,\% worse. By the pre-declared rule CLIA became the v1 model, and the
combination was kept as a pre-declared trade-off variant. Among the factor arms, removing the marking-ROI loss raised absorption (by 0.024); removing
histogram matching lowered both absorption (by 0.039) and the synthetic bone error, but that augmentation targets the appearance of real
radiographs, which synthetic validation does not measure; removing the frequency loss or the rotation changed absorption little.

On the test CTs (Table~\ref{tab:v1_h1}) H1 passes: CLIA absorbs 0.022 less than the recipe arm $[0.013,0.032]$, Holm $p<0.001$, with
non-inferior bone MAE. The lesion-weighted loss also lowers absorption (0.019) and is non-inferior on the test CTs; the combination lowers it most
(0.066) and misses the margin again. The 300k-sample finals of CLIA and recipe, trained later with three seeds each, replicate the contrast.
The Holm adjustment in these tables runs over all seven contrasts of the ablation, which is more conservative than the three named in the plan.
\begin{table}[htbp]
\centering\small
\caption{v1 ablation on synthetic validation CTs (selection data; three seeds each, 200k samples): every arm against the recipe arm.}
\label{tab:v1_ablation}
\fittable{tables/v1_ablation_val.tex}
\end{table}
\begin{table}[htbp]
\centering\small
\caption{First analysis plan, H1: lesion-aware arms against the recipe arm on synthetic LIDC test CTs (three seeds each), and the 300k-sample
finals as a replication.}
\label{tab:v1_h1}
\fittable{tables/v1_h1.tex}
\end{table}

\section{Real-image check and the resulting deviation}
After the ablation the v1 arms were compared with the released weights of the reference model on the 75 SZCH-X-Rays pairs, the comparison
that SAP.md had locked for the final evaluation. The v1 models had the higher bone-region PSNR$_{cal}$ (22.04 against 21.66\,dB) and the lower
full-resolution SSIM$_{cal}$ (0.850 against 0.869), and in the real-image audit on all 75 images they absorbed far more of an inserted lesion
($A$ 0.306 against 0.134 over 450 sites). Because these images had now been looked at, SZCH-X-Rays could no longer serve as a clean locked set
for decisions about a new recipe. It became a development set: 24 images, chosen by a hash of the file name (\path{tools/szch_split.py}),
were used for paired-metric decisions from then on, the other 51 were held out for final fidelity claims, and the audit on all 75 images became
development evidence only. The confirmatory audit moved to VinDr-RibCXR and JSRT, which no model had touched. Appendix~\ref{app:deviations}
lists this deviation with the others.

\section{Consequences for the v2 design}
The v1 study answered its own question: on rendered data with exact truth, lesion insertion lowers absorption a little at no cost in bone
error. It did not solve the problem it was meant to solve. On real radiographs the v1 models absorbed more than twice what the reference model
absorbed, and the absorption hardly depended on the lesion amplitude (Chapter~\ref{ch:audit}): it behaved like a uniform leak of lesion
signal into the bone estimate rather than a failure on hard cases. Three changes followed. The training source moved to the Qure.ai CT
projections, which are translated towards the appearance of real radiographs and are exactly additive (Chapter~\ref{ch:v2}). Lesion insertion,
which had acted only through the supervised target, was replaced by a term that ties the predictions with and without the lesion directly.
And a second analysis plan, SAP\_v2.md, was written before any held-out or locked evaluation of the new models. On the final locked sets the v1
model still absorbs 0.289 of an inserted lesion on VinDr-RibCXR, against 0.149 for the reference model and 0.051 for the v2 ensemble
(Chapter~\ref{ch:audit}).
